\subsection{5.2 精化业务层}
\noindent 业务层按领域边界组织，控制类协调用例，领域对象维护规则，数据映射器负责持久化。订单/合同/逐件交接、库存位置、受损/维修、结算、人员运营分别形成清晰聚合；对象业务不直接调用页面组件或把 SQL 语句当规则。

\begin{BSRTable}
\rowcolor{TableHead}\color{white}\bfseries 领域服务 & \color{white}\bfseries 不变量、命令与事件 \\
预约与报价服务 & `quote(window,store,lines)` 查整个时段可订量；`hold(consent,version,key)` 原子再核价/库存并记暂留到期。修改时段使旧报价失效。\\
合同与履约服务 & `issue/signCurrent` 只签当前版本；`prepare/bind/handover` 检查身份码、签署、安全及单件唯一占用；`returnUnit` 逐件冻结时刻与当前店。\\
鉴定与维修服务 & 前台记录事实，技师在 24h 内初评并附部位照片，经理复核费用；工单冻结费率，技师提交后继续隔离，经理安全验收才恢复可租。\\
结算服务 & 每件超时费+核定损坏费形成总费用；押金抵扣取较小值，补差由专员联系；有争议、未还或回执不明则停在待核，不能自动显示到账。\\
人员与提醒服务 & 班次是计划，服务器打卡是事实；岗位/对象级授权控制读写；提醒已读只影响本人新增标记，业务未结由对象状态判定。\\
\end{BSRTable}

\subsubsection{5.2.1 跨角色一致性与失败恢复}
\noindent 同一服务端事务写入订单、单件、状态事件与审计；各角色工作台通过重新获取授权读模型查看统一事实。页面缓存带版本，旧版本写入拒绝并要求重读。支付、退款、通知和签署属外部副作用，目标采用幂等回执与“待核实/失败”状态恢复；当前本地路径只模拟回执，不能由 UI 文案宣称真实收款。

\subsubsection{5.2.2 政策版本与金额可复核}
\noindent 费用规则与合同条款版本随订单冻结。新客宽限仅首笔\textbf{有效生效租赁}使用，取消或暂留不消耗资格；超时费用逐件按实还时刻计算并封顶于相应日租金，具体百分比待 BSR 批准。受损无损/正常磨损零费，轻微至不可修复以工时、材料与折旧估值为依据，经理核费前不扣押金。旧单没有费率版本时走人工核对。
